\documentclass[11pt,a4paper]{amsart}

\usepackage[T1]{fontenc}
\usepackage[margin=1in]{geometry}
\usepackage{amsmath,amssymb,mathtools}
\usepackage{microtype}
\usepackage[hidelinks]{hyperref}

\newtheorem{theorem}{Theorem}[section]
\newtheorem{proposition}[theorem]{Proposition}
\newtheorem{lemma}[theorem]{Lemma}
\newtheorem{corollary}[theorem]{Corollary}
\theoremstyle{definition}
\newtheorem{definition}[theorem]{Definition}
\newtheorem{remark}[theorem]{Remark}
\newtheorem{example}[theorem]{Example}

\newcommand{\Z}{\mathbb Z}
\newcommand{\N}{\mathbb N}
\newcommand{\Acal}{\mathcal A}
\newcommand{\Fcal}{\mathcal F}
\newcommand{\Scal}{\mathcal S}
\newcommand{\eps}{\varepsilon}
\newcommand{\leg}[2]{\left(\frac{#1}{#2}\right)}

\numberwithin{equation}{section}

\title{Integer Solutions of \texorpdfstring{$z^2-xy^2=x^3-2$}{z2 - xy2 = x3 - 2}}
\author{}
\address{}
\date{}

\subjclass[2020]{Primary 11D09; Secondary 11A15}
\keywords{Diophantine equation, generalized Pell equation, Pell unit, quadratic residue}

\begin{document}

\begin{abstract}
We give an unconditional and self-contained classification of all integer triples satisfying
\(z^2-xy^2=x^3-2\).
Elementary local arguments first force the positive integer \(x\) into a restricted set of residue classes and impose a splitting condition on every prime divisor of \(x\).
For square \(x\), every solution is recovered from a finite divisor calculation; for nonsquare \(x\), every solution is recovered uniquely from one of finitely many reduced seeds followed by a Pell recurrence.
The proof includes an elementary existence proof for positive solutions of Pell's equation, so no external theorem is required.
\end{abstract}

\maketitle

\section{Statement of the classification}

We study
\begin{equation}
 z^2-xy^2=x^3-2,
 \label{eq:main}
\end{equation}
which is equivalently
\begin{equation}
 z^2+2=x(x^2+y^2).
 \label{eq:norm-form}
\end{equation}
The equation is invariant under the independent sign changes \(y\mapsto-y\) and
\(z\mapsto-z\).  It is therefore natural to classify first the normalized solutions
with \(y,z\geq0\), and then restore the signs.

Define the set of admissible positive integers
\begin{equation}
 \Acal=
 \left\{
 x\in\Z_{>0}:
 \begin{array}{l}
 x\text{ is odd},\
 x\not\equiv2\pmod 3,\
 p\mid x\text{ prime}\Longrightarrow p\equiv1\text{ or }3\pmod 8
 \end{array}
 \right\}.
 \label{eq:admissible}
\end{equation}
The term ``admissible'' refers only to necessary local conditions; membership in
\(\Acal\) is not sufficient for solubility.

For an odd integer \(m\geq3\), put
\begin{equation}
 M_m=m^6-2
 \label{eq:Mm}
\end{equation}
and define the finite set
\begin{equation}
 \Fcal_m=
 \left\{
 a\in\Z_{>0}:
 a\mid M_m,\quad
 a\leq \frac{M_m}{a},\quad
 \frac{M_m}{a}\equiv a\pmod{2m}
 \right\}.
 \label{eq:Fm}
\end{equation}
For \(a\in\Fcal_m\), set
\begin{equation}
 Y_{m,a}=\frac{M_m/a-a}{2m},
 \qquad
 Z_{m,a}=\frac{M_m/a+a}{2}.
 \label{eq:squareYZ}
\end{equation}

Now let \(x\geq2\) be nonsquare.  A positive Pell pair for \(x\) means a pair
\((u,v)\in\Z_{>0}^2\) satisfying
\begin{equation}
 u^2-xv^2=1.
 \label{eq:Pell-x}
\end{equation}
As will be proved in Section~\ref{sec:pell}, such pairs exist.  Among them there is a
unique pair \((u_x,v_x)\) for which
\begin{equation}
 \eps_x=u_x+v_x\sqrt{x}
 \label{eq:fundamental-unit}
\end{equation}
assumes its least value.  Put
\begin{equation}
 N_x=x^3-2
 \label{eq:Nx}
\end{equation}
and define the reduced seed set
\begin{equation}
 \Scal_x=
 \left\{
 (r,s)\in\Z_{>0}\times\Z_{\geq0}:
 r^2-xs^2=N_x,\quad s^2<v_x^2N_x
 \right\}.
 \label{eq:Sx}
\end{equation}
This set is finite.  For \((r,s)\in\Scal_x\) and \(n\geq0\), define
\((Z_n,Y_n)=(Z_n(x;r,s),Y_n(x;r,s))\) by
\begin{equation}
 Z_n+Y_n\sqrt{x}=(r+s\sqrt{x})\eps_x^n.
 \label{eq:Pell-param}
\end{equation}
Equivalently,
\begin{equation}
 \begin{aligned}
 Z_0&=r, & Y_0&=s,\\
 Z_{n+1}&=u_xZ_n+xv_xY_n,
 &Y_{n+1}&=v_xZ_n+u_xY_n.
 \end{aligned}
 \label{eq:recurrence}
\end{equation}

\begin{theorem}[Complete classification]
\label{thm:classification}
Every integer solution of \eqref{eq:main} has \(x\in\Acal\), and all solutions are
obtained as follows.

\begin{enumerate}
\item The solutions with \(x=1\) are precisely
\begin{equation}
 (x,y,z)=(1,\pm1,0).
 \label{eq:x1-solutions}
\end{equation}

\item Let \(x=m^2\), where \(m\geq3\) and \(m^2\in\Acal\).  The solutions with this
value of \(x\) are precisely
\begin{equation}
 (x,y,z)=\bigl(m^2,\sigma Y_{m,a},\tau Z_{m,a}\bigr),
 \qquad
 a\in\Fcal_m,
 \quad \sigma,\tau\in\{\pm1\}.
 \label{eq:square-classification}
\end{equation}
If \(\Fcal_m=\varnothing\), then this fiber has no solutions.

\item Let \(x\in\Acal\) be nonsquare.  The solutions with this value of \(x\) are
precisely
\begin{equation}
 (x,y,z)=\bigl(x,\sigma Y_n(x;r,s),\tau Z_n(x;r,s)\bigr),
 \label{eq:nonsquare-classification}
\end{equation}
where \((r,s)\in\Scal_x\), \(n\in\Z_{\geq0}\), and
\(\sigma,\tau\in\{\pm1\}\).  If \(\Scal_x=\varnothing\), then this fiber has no
solutions.
\end{enumerate}

After replacing \(y,z\) by \(|y|,|z|\), the parameter \(a\) in the square case is
unique, and the pair consisting of a reduced seed \((r,s)\) and an exponent \(n\) in
the nonsquare case is unique.
\end{theorem}

The proof occupies the next three sections.

\section{Local restrictions on \texorpdfstring{$x$}{x}}

We begin with the elementary quadratic-residue fact needed below.  Its proof is
included for completeness.

\begin{lemma}[The residue class of primes representing minus two]
\label{lem:minus-two}
Let \(p\) be an odd prime.  The congruence
\begin{equation}
 t^2\equiv-2\pmod p
 \label{eq:minus-two-congruence}
\end{equation}
has a solution if and only if
\begin{equation}
 p\equiv1\text{ or }3\pmod8.
 \label{eq:minus-two-criterion}
\end{equation}
\end{lemma}

\begin{proof}
We give the standard elementary derivation.  For \(a\not\equiv0\pmod p\), write
\(\leg ap\in\{\pm1\}\) for the Legendre symbol.

First, Fermat's congruence
\begin{equation}
 a^{p-1}\equiv1\pmod p
 \label{eq:Fermat}
\end{equation}
follows because multiplication by \(a\) permutes the nonzero residue classes modulo
\(p\): multiplying all classes gives
\(a^{p-1}(p-1)!\equiv(p-1)!\pmod p\), and \((p-1)!\) is invertible modulo \(p\).
Define \(\leg ap=1\) when \(a\) is a nonzero square modulo \(p\), and
\(\leg ap=-1\) otherwise.  The nonzero squares modulo \(p\) are exactly the
\((p-1)/2\) roots of \(T^{(p-1)/2}-1\): every square is a root by
\eqref{eq:Fermat}, the squares are pairwise distinct up to the identification
\(b\sim-b\), and a nonzero polynomial of degree \(d\) over a field has at most
\(d\) roots, by induction from the factor theorem.  If \(a\) is not a square,
then \(a^{(p-1)/2}\not\equiv1\pmod p\), while its square is \(1\) by
\eqref{eq:Fermat}; since \(p\) is odd, it must therefore be congruent to \(-1\).
Consequently
\begin{equation}
 a^{(p-1)/2}\equiv\leg ap\pmod p.
 \label{eq:Euler-criterion}
\end{equation}
This is Euler's criterion.  It also proves the multiplicativity of the Legendre
symbol.

We next prove the form of Gauss's lemma that we need.  Let \(\nu(a)\) be the number
of integers \(j\in\{1,\ldots,(p-1)/2\}\) for which the least positive residue of
\(aj\) is greater than \(p/2\).  Replacing each such residue by its negative, the
least absolute residues of
\(a,2a,\ldots,((p-1)/2)a\) are, up to order,
\(1,2,\ldots,(p-1)/2\).  Indeed, equality up to sign between two of them would give
\(a(i\mp j)\equiv0\pmod p\), impossible for distinct
\(1\leq i,j\leq(p-1)/2\).  Multiplication and cancellation therefore give
\begin{equation}
 a^{(p-1)/2}\equiv(-1)^{\nu(a)}\pmod p.
 \label{eq:Gauss-lemma}
\end{equation}
Together with \eqref{eq:Euler-criterion}, this is Gauss's lemma.

For \(a=-1\), every one of the \((p-1)/2\) residues lies in the upper half, so
\begin{equation}
 \leg{-1}{p}=(-1)^{(p-1)/2}.
 \label{eq:minus-one-symbol}
\end{equation}
For \(a=2\), the residues \(2j\), with \(1\leq j\leq(p-1)/2\), are all less than
\(p\), and \(2j>p/2\) exactly when \(j>p/4\).  Thus
\begin{equation}
 \nu(2)=\frac{p-1}{2}-\left\lfloor\frac p4\right\rfloor.
 \label{eq:nu-two}
\end{equation}
A direct check in the four odd residue classes modulo \(8\) yields
\begin{equation}
 \leg2p=
 \begin{cases}
 1,&p\equiv1,7\pmod8,\\
 -1,&p\equiv3,5\pmod8.
 \end{cases}
 \label{eq:two-symbol}
\end{equation}
Multiplicativity, \eqref{eq:minus-one-symbol}, and \eqref{eq:two-symbol} now give
\begin{equation}
 \leg{-2}{p}=1
 \quad\Longleftrightarrow\quad
 p\equiv1\text{ or }3\pmod8,
\end{equation}
which is equivalent to the asserted solvability criterion.
\end{proof}

\begin{proposition}[Necessary conditions on \(x\)]
\label{prop:local-restrictions}
If \((x,y,z)\in\Z^3\) satisfies \eqref{eq:main}, then \(x\in\Acal\).  In particular,
\begin{equation}
 x\equiv1,3,9,\text{ or }19\pmod{24}.
 \label{eq:x-mod24}
\end{equation}
\end{proposition}

\begin{proof}
Equation \eqref{eq:norm-form} shows first that \(x>0\).  Indeed, its left-hand side
is positive, whereas its right-hand side is zero when \(x=0\) and negative when
\(x<0\).

Let \(p\) be an odd prime divisor of \(x\).  Reduction of
\eqref{eq:norm-form} modulo \(p\) gives
\begin{equation}
 z^2\equiv-2\pmod p.
 \label{eq:p-divides-x}
\end{equation}
By Lemma~\ref{lem:minus-two}, every such \(p\) satisfies
\begin{equation}
 p\equiv1\text{ or }3\pmod8.
 \label{eq:prime-divisor-condition}
\end{equation}

We next prove that \(x\) is odd.  Suppose to the contrary that \(x\) is even.  Then
\eqref{eq:norm-form} implies that \(z\) is even.  Writing \(z=2w\), we have
\begin{equation}
 z^2+2=2(2w^2+1),
\end{equation}
so the exact exponent of \(2\) dividing \(z^2+2\) is one.  Consequently
\eqref{eq:norm-form} forces
\begin{equation}
 x=2u
 \label{eq:x-2u}
\end{equation}
with \(u\) odd, and it also forces \(y\) to be odd: since \(x^2\) is even, the factor
\(x^2+y^2\) must be odd.  Dividing \eqref{eq:norm-form} by \(2\) gives
\begin{equation}
 2w^2+1=u(4u^2+y^2).
 \label{eq:even-x-reduced}
\end{equation}
Because \(u\) and \(y\) are odd,
\begin{equation}
 4u^2+y^2\equiv5\pmod8.
\end{equation}
The left-hand side of \eqref{eq:even-x-reduced} is congruent to \(1\) modulo \(8\)
when \(w\) is even and to \(3\) modulo \(8\) when \(w\) is odd.  Hence
\begin{equation}
 5u\equiv1\text{ or }3\pmod8,
\end{equation}
which gives
\begin{equation}
 u\equiv5\text{ or }7\pmod8.
 \label{eq:u-bad}
\end{equation}
On the other hand, every prime divisor of \(u\) is an odd prime divisor of \(x\), so
it is congruent to \(1\) or \(3\) modulo \(8\) by
\eqref{eq:prime-divisor-condition}.  The set \(\{1,3\}\) is closed under
multiplication modulo \(8\), because \(3^2\equiv1\pmod8\).  It follows that
\(u\equiv1\) or \(3\pmod8\), contradicting \eqref{eq:u-bad}.  Thus \(x\) is odd.

Finally, suppose that \(x\equiv2\pmod3\).  Reducing \eqref{eq:main} modulo \(3\)
gives
\begin{equation}
 z^2+y^2\equiv0\pmod3.
\end{equation}
As the only quadratic residues modulo \(3\) are \(0\) and \(1\), this forces
\(3\mid y\) and \(3\mid z\).  Thus the left-hand side of
\(z^2-xy^2=x^3-2\) is divisible by \(9\).  But every integer congruent to \(2\)
modulo \(3\) has cube congruent to \(8\) modulo \(9\), so
\begin{equation}
 x^3-2\equiv6\pmod9,
\end{equation}
a contradiction.  Therefore \(x\not\equiv2\pmod3\), and all conditions defining
\(\Acal\) hold.

Since \(x\) is odd and all its prime divisors are \(1\) or \(3\) modulo \(8\), one
has \(x\equiv1\) or \(3\pmod8\).  Combining this with
\(x\equiv0\) or \(1\pmod3\) gives exactly the four classes in
\eqref{eq:x-mod24}.
\end{proof}

\section{A finite-seed theorem for generalized Pell equations}
\label{sec:pell}

We now develop, from first principles, the Pell machinery used in the nonsquare
fibers.

\begin{lemma}[Elementary rational approximation]
\label{lem:approximation}
Let \(\alpha\) be irrational.  For every integer \(Q\geq1\), there exist integers
\(p,q\) such that
\begin{equation}
 1\leq q\leq Q,
 \qquad
 |q\alpha-p|<\frac1Q\leq\frac1q.
 \label{eq:Dirichlet}
\end{equation}
As \(Q\) varies, the resulting inequalities hold for infinitely many distinct
positive integers \(q\).
\end{lemma}

\begin{proof}
Consider the \(Q+1\) fractional parts
\begin{equation}
 \{0\alpha\},\{1\alpha\},\ldots,\{Q\alpha\}
\end{equation}
in the \(Q\) half-open intervals
\([j/Q,(j+1)/Q)\), \(0\leq j<Q\).  Two fractional parts lie in the same interval.
Subtracting them gives integers \(p,q\), with \(1\leq q\leq Q\), satisfying
\(|q\alpha-p|<1/Q\).  The inequality \(1/Q\leq1/q\) follows from \(q\leq Q\).

If only finitely many values of \(q\) occurred, then the nonzero numbers
\(|q\alpha-p|\) arising from those finitely many denominators would have a positive
minimum.  This contradicts the bound \(|q\alpha-p|<1/Q\) as \(Q\to\infty\).
\end{proof}

\begin{theorem}[Existence of positive Pell solutions]
\label{thm:Pell-existence}
For every positive nonsquare integer \(D\), there exist \(u,v\in\Z_{>0}\) satisfying
\begin{equation}
 u^2-Dv^2=1.
 \label{eq:Pell-general}
\end{equation}
\end{theorem}

\begin{proof}
Apply Lemma~\ref{lem:approximation} to \(\alpha=\sqrt D\).  Discarding finitely many
pairs if necessary, we may assume \(p,q>0\).  There are infinitely many distinct
pairs \((p,q)\) such that
\begin{equation}
 |q\sqrt D-p|<\frac1q.
 \label{eq:Pell-approx}
\end{equation}
For each pair put
\begin{equation}
 k=p^2-Dq^2.
\end{equation}
The integer \(k\) is nonzero because \(D\) is not a square.  Moreover,
\begin{align}
 |k|
 &=|p-q\sqrt D|\,(p+q\sqrt D)\notag\\
 &<\frac1q\left(2q\sqrt D+\frac1q\right)
 <2\sqrt D+1.
 \label{eq:k-bound}
\end{align}
Thus \(k\) ranges over a finite set of nonzero integers.  One value of \(k\) occurs
for infinitely many pairs \((p,q)\).  Among those pairs, two distinct pairs
\((p_1,q_1)\) and \((p_2,q_2)\) are congruent componentwise modulo \(|k|\), by the
pigeonhole principle.

Define
\begin{equation}
 U=\frac{p_1p_2-Dq_1q_2}{k},
 \qquad
 V=\frac{q_1p_2-p_1q_2}{k}.
 \label{eq:UV-construction}
\end{equation}
Both are integers.  Indeed, modulo \(|k|\),
\begin{equation}
 p_1p_2-Dq_1q_2\equiv p_1^2-Dq_1^2=k\equiv0,
\end{equation}
and
\begin{equation}
 q_1p_2-p_1q_2\equiv q_1p_1-p_1q_1=0.
\end{equation}
The identity
\begin{equation}
 U+V\sqrt D
 =\frac{(p_1+q_1\sqrt D)(p_2-q_2\sqrt D)}{k}
 \label{eq:UV-product}
\end{equation}
then gives, after multiplication by its conjugate,
\begin{equation}
 U^2-DV^2
 =\frac{(p_1^2-Dq_1^2)(p_2^2-Dq_2^2)}{k^2}=1.
 \label{eq:UV-norm}
\end{equation}
We have \(V\neq0\).  Otherwise \(p_1/q_1=p_2/q_2\).  Writing this common positive
rational number in lowest terms as \(a/b\), there would be positive integers
\(t_1,t_2\) with \((p_i,q_i)=t_i(a,b)\).  Since both pairs have the same nonzero
value \(k\),
\begin{equation}
 t_1^2(a^2-Db^2)=t_2^2(a^2-Db^2),
\end{equation}
so \(t_1=t_2\), contrary to the pairs being distinct.  Equation
\eqref{eq:UV-norm} also implies \(U\neq0\).  Therefore
\((u,v)=(|U|,|V|)\) is a positive solution of \eqref{eq:Pell-general}.
\end{proof}

\begin{definition}[Least positive Pell unit]
\label{def:fundamental}
Let \(D\) be a positive nonsquare integer.  Among all positive solutions of
\(u^2-Dv^2=1\), choose the one minimizing \(u+v\sqrt D\), and denote it by
\((u_D,v_D)\).  We write
\begin{equation}
 \eps_D=u_D+v_D\sqrt D.
 \label{eq:epsilon-D}
\end{equation}
This minimum exists: Theorem~\ref{thm:Pell-existence} supplies at least one positive
solution, and below any fixed real bound there are only finitely many integer pairs
\((u,v)\).  Since
\(u_D^2-Dv_D^2=1\), one has
\begin{equation}
 \eps_D^{-1}=u_D-v_D\sqrt D\in\Z[\sqrt D].
 \label{eq:epsilon-inverse}
\end{equation}
\end{definition}

\begin{theorem}[Finite-seed decomposition]
\label{thm:finite-seed}
Let \(D,N\in\Z_{>0}\), with \(D\) nonsquare.  Let
\(\eps_D=u_D+v_D\sqrt D\) be as in Definition~\ref{def:fundamental}, and define
\begin{equation}
 \mathcal R(D,N)=
 \left\{
 (r,s)\in\Z_{>0}\times\Z_{\geq0}:
 r^2-Ds^2=N,\quad s^2<v_D^2N
 \right\}.
 \label{eq:RDN}
\end{equation}
Then \(\mathcal R(D,N)\) is finite.  Every pair
\((Z,Y)\in\Z_{>0}\times\Z_{\geq0}\) satisfying
\begin{equation}
 Z^2-DY^2=N
 \label{eq:generalized-Pell}
\end{equation}
has a unique representation
\begin{equation}
 Z+Y\sqrt D=(r+s\sqrt D)\eps_D^n
 \label{eq:finite-seed-representation}
\end{equation}
with \((r,s)\in\mathcal R(D,N)\) and \(n\in\Z_{\geq0}\).  Conversely, every
right-hand side of \eqref{eq:finite-seed-representation} has nonnegative integral
coefficients and yields a solution of \eqref{eq:generalized-Pell}.
\end{theorem}

\begin{proof}
The set \(\mathcal R(D,N)\) is finite because
\(0\leq s<v_D\sqrt N\), and for each such integer \(s\) the equation
\(r^2=N+Ds^2\) determines at most one positive integer \(r\).

Let \((Z,Y)\) satisfy \eqref{eq:generalized-Pell}, and put
\begin{equation}
 \alpha=Z+Y\sqrt D.
\end{equation}
Since \(N>0\), one has \(Z^2>DY^2\), hence
\(Z>Y\sqrt D\).  Thus both \(\alpha\) and its conjugate are positive, their product
is \(N\), and \(\alpha\geq\sqrt N\).  Because \(\eps_D>1\), there is a unique
integer \(n\geq0\) such that
\begin{equation}
 \eps_D^n\sqrt N\leq\alpha<\eps_D^{n+1}\sqrt N.
 \label{eq:reduction-interval-alpha}
\end{equation}
Set
\begin{equation}
 \gamma=\alpha\eps_D^{-n}.
\end{equation}
By \eqref{eq:epsilon-inverse}, \(\gamma\in\Z[\sqrt D]\), so
\begin{equation}
 \gamma=r+s\sqrt D
 \label{eq:gamma-rs}
\end{equation}
for integers \(r,s\).  Its norm is \(N\), and
\eqref{eq:reduction-interval-alpha} gives
\begin{equation}
 \sqrt N\leq\gamma<\eps_D\sqrt N.
 \label{eq:gamma-interval}
\end{equation}
The conjugate of \(\gamma\) is \(N/\gamma\), so
\begin{equation}
 r=\frac12\left(\gamma+\frac N\gamma\right)>0,
 \qquad
 s=\frac1{2\sqrt D}\left(\gamma-\frac N\gamma\right)\geq0.
 \label{eq:r-s-signs}
\end{equation}

It remains to translate the upper bound in \eqref{eq:gamma-interval}.  The function
\begin{equation}
 f(t)=t-\frac Nt
\end{equation}
is strictly increasing for \(t>0\), because
\(f'(t)=1+N/t^2>0\).  Hence
\begin{align}
 2s\sqrt D
 &=\gamma-\frac N\gamma\notag\\
 &<\eps_D\sqrt N-\frac{N}{\eps_D\sqrt N}\notag\\
 &=\sqrt N\left(\eps_D-\eps_D^{-1}\right)
 =2v_D\sqrt{DN}.
 \label{eq:seed-bound-derived}
\end{align}
Therefore \(s<v_D\sqrt N\), equivalently \(s^2<v_D^2N\), and
\((r,s)\in\mathcal R(D,N)\).  This proves existence of the representation.

Conversely, let \((r,s)\in\mathcal R(D,N)\), and put
\(\gamma=r+s\sqrt D\).  Since \(r^2-Ds^2=N>0\), one has
\(r>s\sqrt D\), so the conjugate of \(\gamma\) is positive.  Moreover
\(s\geq0\), and hence \(\gamma\geq\sqrt N\).  The seed inequality gives
\begin{equation}
 f(\gamma)=2s\sqrt D<2v_D\sqrt{DN}
 =\sqrt N(\eps_D-\eps_D^{-1})=f(\eps_D\sqrt N).
 \label{eq:seed-implies-interval}
\end{equation}
Since \(f\) is strictly increasing, it follows that
\begin{equation}
 \sqrt N\leq\gamma<\eps_D\sqrt N.
 \label{eq:seed-reduction-interval}
\end{equation}
Thus every member of the stated seed set lies in the same half-open reduction
interval used above.  Multiplication by \(\eps_D^n\) gives integral coefficients; more
explicitly, if
\begin{equation}
 Z_n+Y_n\sqrt D=(r+s\sqrt D)\eps_D^n,
\end{equation}
then
\begin{equation}
 \begin{aligned}
 Z_{n+1}&=u_DZ_n+Dv_DY_n,\\
 Y_{n+1}&=v_DZ_n+u_DY_n,
 \end{aligned}
\end{equation}
so induction gives \(Z_n>0\) and \(Y_n\geq0\).  Norm multiplicativity gives
\begin{equation}
 Z_n^2-DY_n^2=(r^2-Ds^2)(u_D^2-Dv_D^2)^n=N.
\end{equation}

For uniqueness, \eqref{eq:seed-reduction-interval} shows that any representation
\(\alpha=(r+s\sqrt D)\eps_D^n\) with a stated seed satisfies
\begin{equation}
 \eps_D^n\sqrt N\leq\alpha<\eps_D^{n+1}\sqrt N.
\end{equation}
Hence \(n\) is the unique integer determined by
\eqref{eq:reduction-interval-alpha}.  Once \(n\) is fixed,
\(r+s\sqrt D=\alpha\eps_D^{-n}\) is fixed.  Thus both the seed and the exponent are
unique.
\end{proof}

\begin{corollary}
\label{cor:fiber-infinite}
For fixed positive nonsquare \(D\) and positive \(N\), the equation
\(Z^2-DY^2=N\) has either no solution in
\(\Z_{>0}\times\Z_{\geq0}\) or infinitely many such solutions.
\end{corollary}

\begin{proof}
By Theorem~\ref{thm:finite-seed}, a solution exists if and only if the finite seed
set is nonempty.  From any seed, the quantities
\((r+s\sqrt D)\eps_D^n\) are strictly increasing as \(n\) increases, and therefore
give infinitely many distinct solutions.
\end{proof}

\section{Proof of the classification}

\begin{proof}[Proof of Theorem~\ref{thm:classification}]
Proposition~\ref{prop:local-restrictions} proves that every solution has
\(x\in\Acal\).  It remains to classify each possible fiber.

\smallskip
\noindent\emph{The fiber \(x=1\).}
Equation \eqref{eq:main} becomes
\begin{equation}
 z^2-y^2=-1,
\end{equation}
so
\begin{equation}
 (z-y)(z+y)=-1.
\end{equation}
The only ordered factor pairs of \(-1\) are \((1,-1)\) and \((-1,1)\).  Solving
gives \(z=0\) and \(y=\pm1\), proving \eqref{eq:x1-solutions}.

\smallskip
\noindent\emph{Square fibers.}
Suppose \(x=m^2\), with \(m\geq3\), and write
\begin{equation}
 Y=|y|,
 \qquad
 Z=|z|.
\end{equation}
Then
\begin{equation}
 Z^2-m^2Y^2=m^6-2=M_m.
 \label{eq:square-factor-equation}
\end{equation}
Since \(M_m>0\), one has \(Z>mY\).  Thus
\begin{equation}
 a=Z-mY,
 \qquad
 b=Z+mY
 \label{eq:a-b-definition}
\end{equation}
are positive integers satisfying
\begin{equation}
 ab=M_m,
 \qquad
 a\leq b,
 \qquad
 b-a=2mY.
 \label{eq:a-b-properties}
\end{equation}
Therefore \(b=M_m/a\) and
\begin{equation}
 \frac{M_m}{a}\equiv a\pmod{2m},
\end{equation}
so \(a\in\Fcal_m\).  Solving \eqref{eq:a-b-definition} for \(Y,Z\) yields exactly
\eqref{eq:squareYZ}.

Conversely, let \(a\in\Fcal_m\), and put \(b=M_m/a\).  Because \(m\) is odd,
\(M_m\) is odd; hence its positive divisors \(a,b\) are odd.  Thus
\begin{equation}
 Z=\frac{a+b}{2}\in\Z.
\end{equation}
The congruence in \eqref{eq:Fm} gives
\begin{equation}
 Y=\frac{b-a}{2m}\in\Z_{\geq0}.
\end{equation}
Moreover,
\begin{equation}
 Z-mY=a,
 \qquad
 Z+mY=b,
\end{equation}
and therefore
\begin{equation}
 Z^2-m^2Y^2=ab=M_m.
\end{equation}
This proves the square-fiber parameterization and its completeness.  The normalized
factor \(a=Z-mY\) is determined uniquely by \((Y,Z)\).

\smallskip
\noindent\emph{Nonsquare fibers.}
Let \(x\in\Acal\) be nonsquare.  Since \(x\geq2\), the integer
\(N_x=x^3-2\) is positive.  With \(Y=|y|\) and \(Z=|z|\), equation
\eqref{eq:main} becomes
\begin{equation}
 Z^2-xY^2=N_x.
 \label{eq:fiber-generalized-Pell}
\end{equation}
Apply Theorem~\ref{thm:finite-seed} with \(D=x\) and \(N=N_x\).  Its seed set is
precisely \(\Scal_x\), and its representation is precisely
\eqref{eq:Pell-param}, equivalently \eqref{eq:recurrence}.  It gives every normalized
solution and proves the uniqueness of the seed and exponent.  Restoring the
independent signs of \(y\) and \(z\) completes the proof.
\end{proof}

\begin{corollary}[Structure of the fibers]
\label{cor:fibers}
Every square fiber of \eqref{eq:main} is finite.  Every nonsquare fiber is either
empty or infinite.  The full equation has infinitely many integer solutions.
\end{corollary}

\begin{proof}
The square assertion follows from the finiteness of the divisor set
\(\Fcal_m\).  The nonsquare assertion follows from
Corollary~\ref{cor:fiber-infinite}.  Infinitely many solutions occur already in the
fiber \(x=3\), as shown explicitly below.
\end{proof}

\section{Examples and effective use of the theorem}

The classification is effective for each fixed \(x\).  In a square fiber, one
factors \(m^6-2\) and checks the finite congruence in \eqref{eq:Fm}.  In a nonsquare
fiber, one searches for the least positive solution of \(u^2-xv^2=1\); existence is
guaranteed by Theorem~\ref{thm:Pell-existence}.  One then checks the finitely many
integers \(s\geq0\) satisfying \(s^2<v_x^2N_x\), retaining those for which
\(N_x+xs^2\) is a square.

\begin{example}[The complete fiber \(x=3\)]
Here
\begin{equation}
 N_3=25,
 \qquad
 \eps_3=2+\sqrt3.
\end{equation}
Indeed, \(2^2-3\cdot1^2=1\), and no smaller positive Pell pair is possible because
\(v\geq1\) and \(u^2=1+3v^2\geq4\).
A reduced seed satisfies
\begin{equation}
 r^2-3s^2=25,
 \qquad
 0\leq s<5.
 \label{eq:x3-seed}
\end{equation}
Modulo \(5\), this gives \(r^2\equiv3s^2\pmod5\).  Since the nonzero squares modulo
\(5\) are \(1\) and \(4\), the number \(3\) is a nonsquare; hence \(5\mid s\), and
then \(5\mid r\).  The bound in \eqref{eq:x3-seed} forces \(s=0\), and then
\(r=5\).  Thus all normalized solutions with \(x=3\) are
\begin{equation}
 Z_n+Y_n\sqrt3=5(2+\sqrt3)^n,
 \qquad n\geq0.
 \label{eq:x3-family}
\end{equation}
The first pairs \((Y_n,Z_n)\) are
\begin{equation}
 (0,5),\ (5,10),\ (20,35),\ (75,130),\ (280,485),\ldots.
\end{equation}
Independent signs give all integer solutions in this fiber.
\end{example}

\begin{example}[The complete fiber \(x=9\)]
Here \(m=3\) and
\begin{equation}
 M_3=3^6-2=727.
\end{equation}
The integer \(727\) has no prime divisor at most \(\sqrt{727}<27\): it is not
divisible by \(2,3,5,7,11,13,17,19\), or \(23\).  Hence it is prime.  The only
normalized factor pair is \((a,b)=(1,727)\), and
\(727\equiv1\pmod6\).  Formula \eqref{eq:squareYZ} gives
\begin{equation}
 Y=\frac{727-1}{6}=121,
 \qquad
 Z=\frac{727+1}{2}=364.
\end{equation}
Therefore all solutions with \(x=9\) are
\begin{equation}
 (9,\pm121,\pm364).
\end{equation}
\end{example}

\begin{example}[Admissibility is not sufficient]
The prime \(43\) belongs to \(\Acal\), because
\(43\equiv3\pmod8\) and \(43\equiv1\pmod3\).  Nevertheless the fiber \(x=43\) is
empty.  Indeed,
\begin{equation}
 43^3-2=79505\equiv5\pmod{25}.
\end{equation}
If a solution existed, reduction modulo \(5\) would give
\begin{equation}
 z^2-3y^2\equiv0\pmod5.
\end{equation}
As \(3\) is a nonsquare modulo \(5\), this forces \(5\mid y\) and \(5\mid z\).
The left-hand side \(z^2-43y^2\) would then be divisible by \(25\), whereas the
right-hand side is congruent to \(5\pmod{25}\), a contradiction.
\end{example}

\end{document}
